\chapter{EXPERIMENTAL WORK AND IMPLEMENTATION}
\label{workdone}
\subsection{Preliminary Classification and ROI Extraction}
The first classification experiments were trained on raw B-mode ultrasound images. These results suggested that performance could improve if the model focused on the LMM rather than the entire ultrasound frame. This led to the ROI extraction stage.

A U-Net segmentation model with a ResNet-34 encoder pretrained on ImageNet was trained on LUMINOUS images with annotated GMC images. The images were resized to 256 $\times$ 256 pixels and trained with a batch size of 8. On LUMINOUS, the model achieved IoU = 0.86, Dice = 0.92, and Pixel Accuracy = 0.98. However, when tested on GMC images, the segmentation quality dropped noticeably. The predicted masks were noisy in several cases, showing a clear domain shift between the public LUMINOUS data and the local clinical images.

\begin{figure}[H]
\centering
\safeincludegraphics[width=0.2\textwidth]{bmode.png}
\caption{B-mode USG image}
\label{fig:bmode}
\end{figure}

\begin{figure}[H]
\centering
\safeincludegraphics[width=0.2\textwidth]{groundtruthmask.png}
\caption{Ground truth mask}
\label{fig:groundtruth}
\end{figure}

\begin{figure}[H]
\centering
\safeincludegraphics[width=0.2\textwidth]{predicted mask.png}
\caption{Predicted mask}
\label{fig:predmask}
\end{figure}

\begin{figure}[H]
\centering
\safeincludegraphics[width=0.2\textwidth]{roi_cut.jpg}
\caption{ROI crop}
\label{fig:roi}
\end{figure}

\subsection{ROI Extraction: Model Progression}
\subsubsection{SAM}
As an alternative to U-Net, the Segment Anything Model (SAM) was tested using bounding box prompts generated from LUMINOUS segmentation masks. SAM achieved strong performance on the test set, with IoU = 0.91, Dice = 0.95, and Pixel Accuracy = 0.99. The main limitation was that SAM needed a bounding box prompt during inference. This meant that another model would still be required to generate the prompt for unseen images. Like U-Net, SAM also struggled with the domain difference between LUMINOUS and GMC images, so it was not suitable as a standalone deployment model.

\subsubsection{YOLO-based ROI Extraction}
Since the final requirement was a reliable muscle crop rather than a pixel-perfect mask, YOLOv8 object detection was selected for ROI extraction. Bounding box annotations were automatically generated from the LUMINOUS masks and used to fine-tune a pretrained YOLOv8m model. The dataset was split 80/20 for training and validation. Images were converted from 16-bit to 8-bit and transformed into 3-channel inputs. The model was trained for 100 epochs with an input resolution of 640 $\times$ 640 and batch size 16.

YOLOv8 reached mAP50 = 0.995, Precision = 0.999, and Recall = 1.0 on the validation set. More importantly, it predicted bounding boxes directly during inference, making it more practical for the final pipeline than SAM.

\subsection{Binary Classification}
Three pretrained backbone encoders were evaluated for binary classification, each followed by a three-layer fully connected head. EfficientNet-B0 and ViT-B16 both reached 90.47\% accuracy. ViT-B16 produced slightly better precision, recall, and F1-score values, with all three around 0.94. SigLIP reached 85.71\% accuracy.

A short self-supervised inpainting experiment was also carried out on LUMINOUS and GMC images to reduce the gap between ImageNet-pretrained weights and ultrasound images. Because the available pretraining data was limited, the result was modest, with 81.58\% accuracy and 0.87 F1-score. This direction was not continued in the main pipeline.

\subsection{Heckmatt Grade Classification}
The classification task was then extended to four Heckmatt grades. The dataset contained 606 images, distributed as Grade I: 146, Grade II: 279, Grade III: 129, and Grade IV: 52. This imbalance was a major difficulty, especially for Grade IV.

\subsubsection{Earlier Backbone Experiments}
The first four-class experiments compared several pretrained backbones. Table~\ref{tab:heckmatt_old_results} shows the earlier results.

\begin{table}[H]
\centering
\\resizebox{\textwidth}{!}{%
\begin{tabular}{lcc}
\toprule
\textbf{Model} & \textbf{Best Val Accuracy} & \textbf{Test Accuracy} \\
\midrule
ViT-B16 (cropped, LUMINOUS) & 64.00\% & 47.00\% \\
MedGemma (vision encoder) & 62.00\% & 50.00\% \\
ViT-B16 (full images) & 61.11\% & 55.00\% \\
ResNet & 53.30\% & 51.10\% \\
EfficientNet-B0 & 52.00\% & 50.00\% \\
\bottomrule
\end{tabular}}
\caption{Earlier Heckmatt grading results across backbone architectures}
\label{tab:heckmatt_old_results}
\end{table}

ViT-B16 trained on full images gave the best test accuracy among these initial runs. The cropped ViT-B16 run had the highest validation accuracy but dropped sharply on the test set, which indicated overfitting. MedGemma did not improve the result despite medical pretraining, suggesting that domain-specific pretraining alone was not enough when the fine-tuning dataset was small.

\subsubsection{Direct Four-Class Classification}
The four-class classifiers were first checked in the usual way, using the test split and confusion matrices. On the equalized dataset, the multiclass model reached 50.55\% test accuracy. A ResNet18 direct four-class model gave the strongest direct result, with 63.74\% test accuracy and 67.03\% best validation accuracy. These numbers were useful, but the confusion matrices were more informative than the single accuracy value. H4 was being identified quite reliably even though it had the fewest examples, while most errors were concentrated around H1/H2 and H2/H3.

\begin{table}[H]
\centering
\\resizebox{\textwidth}{!}{%
\begin{tabular}{lcc}
\toprule
\textbf{Experiment} & \textbf{Test Accuracy} & \textbf{Test Loss} \\
\midrule
Equalized multiclass test (vit) & 50.55\% & 2.4477 \\
ResNet18 direct four-class & 63.74\% & 1.6417 \\
Augmentation run (ViT) & 60.20\% & -- \\
ViT + gaussian blur & 57.14\% & 2.102 \\
\bottomrule
\end{tabular}}
\caption{Direct multiclass classification results}
\label{tab:direct_multiclass}
\end{table}

\begin{table}[H]
\centering
\scriptsize
\\resizebox{\textwidth}{!}{%
\begin{tabular}{lcccc}
\toprule
 & \textbf{Pred H1} & \textbf{Pred H2} & \textbf{Pred H3} & \textbf{Pred H4} \\
\midrule
\textbf{True H1} & 11 & 7 & 4 & 0 \\
\textbf{True H2} & 13 & 17 & 11 & 1 \\
\textbf{True H3} & 1 & 4 & 10 & 4 \\
\textbf{True H4} & 0 & 0 & 0 & 8 \\
\bottomrule
\end{tabular}}
\caption{Equalized dataset multiclass confusion matrix}
\label{tab:equalized_cm}
\end{table}

\begin{table}[H]
\centering
\scriptsize
\\resizebox{\textwidth}{!}{%
\begin{tabular}{lcccc}
\toprule
 & \textbf{Pred H1} & \textbf{Pred H2} & \textbf{Pred H3} & \textbf{Pred H4} \\
\midrule
\textbf{True H1} & 12 & 10 & 0 & 0 \\
\textbf{True H2} & 5 & 33 & 3 & 1 \\
\textbf{True H3} & 1 & 8 & 6 & 4 \\
\textbf{True H4} & 0 & 1 & 0 & 7 \\
\bottomrule
\end{tabular}}
\caption{ResNet18 direct four-class confusion matrix}
\label{tab:resnet_direct_cm}
\end{table}

\subsection{Hyperparameter Search}
A 36-run hyperparameter search was carried out using a ViT classification pipeline with ImageNet normalization and Gaussian blur preprocessing. Learning rate, dropout, batch size, and the number of unfrozen final blocks were varied. Trial 32 gave the best validation accuracy of 68.10\%, but it did not carry the same advantage to the test set. Trials 15 and 19 were more stable on the test split, both reaching 61.50\% accuracy and 65.30\% balanced accuracy. Gaussian blur with a radius of 0.25 showed good results.



\begin{table}[H]
\centering
\scriptsize
\\resizebox{\textwidth}{!}{%
\begin{tabular}{lcccccc}
\toprule
\textbf{Trial} & \textbf{LR} & \textbf{Dropout} & \textbf{Blocks} & \textbf{Batch} & \textbf{Test Acc.} & \textbf{Bal. Acc.} \\
\midrule
15 & 0.0003 & 0.2 & 2 & 8 & 61.50\% & 65.30\% \\
19 & 0.0003 & 0.3 & 2 & 8 & 61.50\% & 65.30\% \\
29 & 0.0010 & 0.3 & 1 & 8 & 60.40\% & 65.10\% \\
25 & 0.0010 & 0.2 & 1 & 8 & 60.40\% & 63.10\% \\
32 & 0.0010 & 0.3 & 2 & 16 & 54.90\% & 55.90\% \\
\bottomrule
\end{tabular}}
\caption{Selected hyperparameter search results}
\label{tab:hparam_search}
\end{table}

\subsection{Pairwise Binary Experiments}
Pairwise classifiers were trained to check whether the difficulty came from the model setup or from the grade boundaries themselves. The ViT-B/16 pairwise runs made the pattern clearer. H1 vs H4 and H2 vs H4 were the easiest comparisons, both reaching 93.75\% test accuracy. The weak point was H2 vs H3, which reached only 60.53\%. This supports the visual observation that the middle Heckmatt grades do not always have a clean separation in the available images.

\begin{table}[H]
\centering
\scriptsize
\\resizebox{\textwidth}{!}{%
\begin{tabular}{lcccc}
\toprule
\textbf{Pair} & \textbf{Accuracy} & \textbf{Bal. Acc.} & \textbf{Sensitivity} & \textbf{Specificity} \\
\midrule
H1 vs H2 & 70.45\% & 70.45\% & 81.80\% & 59.10\% \\
H1 vs H3 & 84.21\% & 84.21\% & 78.90\% & 89.50\% \\
H1 vs H4 & 93.75\% & 93.70\% & 100.00\% & 87.50\% \\
H2 vs H3 & 60.53\% & 60.53\% & 78.90\% & 42.10\% \\
H2 vs H4 & 93.75\% & 93.70\% & 87.50\% & 100.00\% \\
H3 vs H4 & 62.50\% & 62.50\% & 87.50\% & 37.50\% \\
\bottomrule
\end{tabular}}
\caption{ViT-B/16 pairwise binary classification results}
\label{tab:vit_pairwise}
\end{table}

\begin{figure}[H]
    \centering
    \safeincludegraphics[width=0.95\textwidth]{heatmap_pairwise.png}
    \caption{Pairwise Heatmap}
    \label{fig:diagrams}
\end{figure}

\subsection{Hierarchical Inference}
Hierarchical inference was tested because Heckmatt grades are ordered rather than independent categories. The low-to-high cascade first checks H1, then H2, then H3/H4. The high-to-low cascade first checks H4, then H3, then H2/H1. This design is shown in Figure~\ref{fig:hierarchical_arch}. The high-to-low ResNet18 cascade performed better overall, reaching 61.54\% accuracy compared with 53.85\% for the low-to-high ResNet18 cascade and 52.75\% for the earlier EfficientNet cascade. However, the cascade logs also showed that the middle decisions were unstable, especially where the model had to choose between adjacent grades.

\begin{table}[H]
\centering
\scriptsize
\\resizebox{\textwidth}{!}{%
\begin{tabular}{lccccc}
\toprule
\textbf{Cascade} & \textbf{Accuracy} & \textbf{H1 Recall} & \textbf{H2 Recall} & \textbf{H3 Recall} & \textbf{H4 Recall} \\
\midrule
EfficientNet low-to-high & 52.75\% & 13.60\% & 73.80\% & 42.10\% & 75.00\% \\
ResNet18 low-to-high & 53.85\% & 54.50\% & 57.10\% & 26.30\% & 100.00\% \\
ResNet18 high-to-low & 61.54\% & 50.00\% & 71.40\% & 36.80\% & 100.00\% \\
\bottomrule
\end{tabular}}
\caption{Hierarchical cascade comparison}
\label{tab:cascade_comparison}
\end{table}

\begin{table}[H]
\centering
\scriptsize
\\resizebox{\textwidth}{!}{%
\begin{tabular}{lcccc}
\toprule
 & \textbf{Pred H1} & \textbf{Pred H2} & \textbf{Pred H3} & \textbf{Pred H4} \\
\midrule
\textbf{True H1} & 11 & 11 & 0 & 0 \\
\textbf{True H2} & 9 & 30 & 2 & 1 \\
\textbf{True H3} & 4 & 5 & 7 & 3 \\
\textbf{True H4} & 0 & 0 & 0 & 8 \\
\bottomrule
\end{tabular}}
\caption{ResNet18 high-to-low cascade confusion matrix}
\label{tab:resnet_high_low_cm}
\end{table}

\subsection{Label Audit}
The label audit was added after the classification experiments showed the same kinds of mistakes repeatedly. The aim was not to claim that the model was always correct, but to find images where the current labels and the model behaviour strongly disagreed. The workflow is shown in Figure~\ref{fig:audit_arch}. H1 and H2 were treated as normal, while H3 and H4 were treated as abnormal. The audit covered 606 total images across GMC and LUMINOUS. It found 172 disagreements, giving a combined disagreement rate of 28.38\%. Of these, 79 were high-confidence disagreements with confidence at least 0.8, giving a high-confidence disagreement rate of 13.04\%.

\begin{table}[H]
\centering
\scriptsize
\resizebox{0.9\textwidth}{!}{%
\begin{tabular}{lc}
\toprule
\textbf{Metric} & \textbf{Value} \\
\midrule
Total images & 606 \\
Disagreements & 172 \\
Disagreement rate & 28.38\% \\
High-confidence disagreements & 79 \\
High-confidence disagreement rate & 13.04\% \\
\bottomrule
\end{tabular}}
\caption{Combined GMC and LUMINOUS image audit summary}
\label{tab:audit_summary}
\end{table}

\begin{table}[H]
\centering
\scriptsize
\\resizebox{\textwidth}{!}{%
\begin{tabular}{lccccc}
\toprule
\textbf{Set/Class} & \textbf{Assumed} & \textbf{Total} & \textbf{Agree} & \textbf{Disagree} & \textbf{High-conf.} \\
\midrule
GMC H1 & Normal & 64 & 40 & 24 & 6 \\
GMC H2 & Normal & 77 & 19 & 58 & 34 \\
GMC H3 & Abnormal & 54 & 46 & 8 & 2 \\
GMC H4 & Abnormal & 23 & 23 & 0 & 0 \\
LUMINOUS H1 & Normal & 76 & 76 & 0 & 0 \\
LUMINOUS H2 & Normal & 189 & 182 & 7 & 0 \\
LUMINOUS H3 & Abnormal & 56 & 3 & 53 & 27 \\
LUMINOUS H4 & Abnormal & 11 & 2 & 9 & 3 \\
\bottomrule
\end{tabular}}
\caption{Class-wise label audit results}
\label{tab:audit_classwise}
\end{table}

The audit highlights different failure modes in the two datasets. In GMC, H2 produced the largest issue, with 58 disagreements out of 77 images and 34 high-confidence disagreements. In LUMINOUS, H3 produced the largest issue, with 53 disagreements out of 56 images and 27 high-confidence disagreements. GMC H4 and LUMINOUS H1 showed complete agreement with the assumed labels. The clinical reviewer has also indicated that some images are currently being checked and corrected, so these audit results should be read as evidence for label review rather than as a final clinical judgement. The behaviour of H4 is especially important: even with fewer samples, it was graded correctly in several experiments, which suggests that the main issue may lie more in label consistency and boundary cases than in data quantity alone.

\subsection{User Interface}
A clinical user interface was developed to make the pipeline easier to use for physicians and domain experts. The interface supports direct upload of TIFF and BMP ultrasound images, displays the YOLOv8-detected ROI over the input image, and shows the predicted Heckmatt grade. This provides an end-to-end workflow for non-technical users.

Another tool that came out of the project is a data annotation interface for preparing and correcting the dataset. In this tool, ultrasound images can be uploaded and then shown to the doctor one by one for review. The doctor can assign or correct the class label and draw the bounding box around the required muscle region. After annotation, the labels and bounding-box information can be downloaded directly using a button. This is useful for the current correction phase because the reviewer can go through the images in a controlled order and generate updated labels without manually editing separate files.

\begin{figure}[H]
    \centering
    \safeincludegraphics[width=0.95\textwidth]{UI.jpg}
    \caption{UI screenshot}
    \label{fig:ui}
\end{figure}

\begin{figure}[H]
    \centering
    \safeincludegraphics[width=0.95\textwidth]{annotation.png}
    \caption{Annotator Screenshot}
    \label{fig:annotation_ui}
\end{figure}
